\subsection{换基域}

设 \(\mathrm{L}\) 是域 \(\mathrm{K}\) 的一个扩张。设 \(\mathrm{A}\) 与 \(\mathrm{B}\) 是有限次数的 K-代数；则 L-代数 \(\mathrm{A}_{(\mathrm{L})}\) 与 \(\mathrm{B}_{(\mathrm{L})}\) 是有限次数的。L-代数 \(\mathrm{A}_{(\mathrm{L})} \otimes_{\mathrm{L}} \mathrm{B}_{(\mathrm{L})}\) 与 \((\mathrm{A} \otimes_{\mathrm{K}} \mathrm{B})_{(\mathrm{L})}\) 同构（III, §4, No.~1, 第 462 页，命题 3）。L-代数 \(\mathrm{K}_{(\mathrm{L})}\) 同构于 \(\mathrm{L}\)。因此，存在唯一的幺半群同态 \(\rho_{\mathrm{L/K}}\)（从 \(\mathcal{C}_{\mathrm{K}}\) 到 \(\mathcal{C}_{\mathrm{L}}\) 的），使得

\[
\rho_ {\mathrm{L/K}} (\operatorname{cl} (\mathrm{A})) = \operatorname{cl} (\mathrm{A} _ {(\mathrm{L})})\tag{3}
\]

对每个有限次数的 K-代数 A。

若 K-代数 A 与 B 是森田等价的，则 L-代数 \(A_{(L)}\) 与 \(B_{(L)}\) 亦然（VIII, 第 111 页，命题 13, e)）。若 K-代数 A 是中心单的且有限次数的，则 L-代数 \(A_{(L)}\) 亦然（VIII, 第 251 页，注 2）。于是我们可以由 \(\rho_{L/K}\) 得出一个群同态 \(r_{L/K}\)（从 \(\mathrm{Br}(K)\) 到 \(\mathrm{Br}(L)\) 的），使得

\[
r _ {\mathrm{L/K}} ([ \mathrm{A} ]) = [ \mathrm{A} _ {(\mathrm{L})} ]\tag{4}
\]

对每个有限次数的中心单 K-代数 A。

设 \(M\) 是域 \(L\) 的一个扩张。由于标量扩张是传递的（III, §1, No.~5, 第 434 页），我们有关系

\[
r _ {\mathrm{M/K}} = r _ {\mathrm{M/L}} \circ r _ {\mathrm{L/K}}.\tag{5}
\]

设 A 是一个有限次数的中心单 K-代数，L 是 K 的一个扩张。我们说 L 是 A 的一个分裂域（或分裂 A），如果 L-代数 \(\mathrm{A}_{(\mathrm{L})}\) 同构于一个矩阵代数 \(\mathbf{M}_{n}(\mathrm{L})\)（对一个整数 \(n \geqslant 1\)）。用前面所述的记号，这相当于说 A 在 \(\mathrm{Br}(\mathrm{K})\) 中的类属于同态 \(r_{\mathrm{L/K}} : \mathrm{Br}(\mathrm{K}) \to \mathrm{Br}(\mathrm{L})\) 的核。

若 B 相似于 A，则 L 是 A 的分裂域当且仅当它是 B 的分裂域。

若 L 是 A 的分裂域，则 L 的每个扩张都是 A 的分裂域。由 VIII, 第 252 页的定理 1，存在 K 的一个有限次数的伽罗瓦扩张，它是 A 的分裂域，且 K 的每个可分闭包都是 A 的分裂域。

命题 5. —— 设 A 是一个有限次数的中心单 K-代数，L 是 K 的一个有限次数的扩张。下列性质等价：

\begin{enumerate}
\def\labelenumi{(\roman{enumi})}
\item
  扩张 L 是 A 的一个分裂域。
\item
  存在一个有限次数的中心单 K-代数，它相似于 A 且包含一个同构于 L 的极大交换子代数。
\end{enumerate}

我们来证明 (ii) 蕴含 (i)；只需考虑 L 是 A 的一个极大交换子代数的情形。设 \(\psi: A \otimes_{K} A^{\circ} \to \operatorname{End}_{K}(A)\) 是典范同构，它把 \(a \otimes a'\) 映到 K-线性映射 \(x \mapsto axa'\)（VIII, 第 252 页，定理 1）。我们把 A 看作 L 上一个右向量空间；则 \(\psi\) 把 \(A \otimes_{K} L\) 映入子空间 \(\operatorname{End}_{L}(A)\)（\(\operatorname{End}_{K}(A)\) 的），并由限制诱导出 L-代数的一个单射同态 \(\psi': A \otimes_{K} L \to \operatorname{End}_{L}(A)\)。设 n = {[}L : K{]}。由 VIII, 第 262 页的命题 3，我们有 {[}A : L{]} = n，因此 \([A \otimes_{K} L : L] = [A : K] = n^{2}\) 与 \([\operatorname{End}_{L}(A) : L] = n^{2}\)。于是 \(\psi'\) 是一个同构，这就证明了 (i)。

其逆由下面的引理 3 得出。

引理 3. —— 设 A 是一个有限次数的中心单 K-代数，L 是 K 的一个有限次数的扩张且是 A 的一个分裂域。设 V 是一个单 \(\mathrm{A}_{(\mathrm{L})}\)-模，于是自然态射 \(\varphi: \mathrm{A}_{(\mathrm{L})} \to \mathrm{End}_{\mathrm{L}}(\mathrm{V})\) 是一个同构。设 C 是环 \(\mathrm{End}_{\mathrm{A}}(\mathrm{V})\)。则 C 相似于 \(A^{o}\)，且 \(L \otimes 1 \subset A_{(L)}\) 的像是 C 的一个极大交换子代数。

我们把 A 与 \(\mathrm{A}_{(\mathrm{L})}\) 的一个子环等同。我们把 V 看作一个 K-向量空间。环 C 是 \(\varphi(\mathrm{A})\) 在 \(\operatorname{End}_{\mathrm{K}}(\mathrm{V})\) 中的交换子。它是一个有限次数的中心单 K-代数，且同态 \(a \otimes c \mapsto ac\)（从 \(A \otimes_{K} C\) 到 \(\operatorname{End}_{\mathrm{K}}(\mathrm{V})\) 的）是一个同构（VIII, 第 260 页，定理 6, a)）。因此，K-代数 A 与 \(C^{o}\) 是相似的（VIII, 第 279 页，命题 3）。

设 \(L_{V}\) 是 L-向量空间 V 的位似环；它是 \(\mathrm{End}_{\mathrm{L}}(\mathrm{V})\) 在 \(\mathrm{End}_{\mathrm{K}}(\mathrm{V})\) 中的交换子（VIII, 第 82 页，推论 1）。而 K-代数 \(\mathrm{End}_{\mathrm{L}}(\mathrm{V})\) 由 \(\varphi(\mathrm{A})\) 与 \(L_{V}\) 生成；因此，在 \(\mathrm{End}_{\mathrm{K}}(\mathrm{V})\) 中，我们有

\[
\mathrm{L} _ {\mathrm{V}} = \mathrm{End} _ {\mathrm{L}} (\mathrm{V}) ^ {\prime} = \varphi (\mathrm{A}) ^ {\prime} \cap \mathrm{L} _ {\mathrm{V}} ^ {\prime} = \mathrm{C} \cap \mathrm{L} _ {\mathrm{V}} ^ {\prime},
\]

其中对 \(\operatorname{End}_{\mathrm{K}}(\mathrm{V})\) 的任何子集 B，B 在 \(\operatorname{End}_{\mathrm{K}}(\mathrm{V})\) 中的交换子记作 \(B'\)。于是 \(L_{V}\) 是 C 的一个极大交换子代数（VIII, 第 261 页，引理 3），因而也是 \(C^{o}\) 的极大交换子代数。映射 \(\lambda \mapsto \lambda_{V}\) 是从 L 到 \(L_{V}\) 的一个 K-代数同构。这就证明了引理。

推论 1. —— 设 A 是一个有限次数的中心单 K-代数，L 是 K 的一个有限次数的扩张。假设 \([A:K]=[L:K]^{2}\)。则 L 是 A 的分裂域当且仅当 A 包含一个同构于 L 的子代数。

假设存在一个从 L 到 A 的态射 \(\varphi\)。设 M 是一个包含 \(\varphi(\mathrm{L})\) 的极大半单交换子代数。由

VIII, 第 262 页的命题 3，我们有 \([A:K] = [M:K]^2\)，因此 \([M:K] = [L:K]\) 且 \(M = \varphi(L)\)。由命题 5，域 \(L\) 分裂 \(A\)。反之，假设域 \(L\) 分裂 \(A\)；则它同构于一个中心单 K-代数 \(B\)（相似于 \(A\) 的）的一个极大交换子代数（命题 5）。我们有 \([B:K] = [L:K]^2\)（VIII, 第 262 页，命题 3），因此 \([B:K] = [A:K]\)。于是 \(B\) 同构于 \(A\)（VIII, 第 280 页，推论）。由此即得推论 1。

推论 2. —— 设 D 是 K 上一个有限次数且以 K 为中心的体，L 是 K 的一个有限次数的扩张且是 D 的一个分裂域。D 的约化次数整除 {[}L : K{]}。

用 r 表示 D 的约化次数（VIII, 第 253 页）；按定义，我们有 \([D:K]=r^{2}\)。由命题 5，存在一个相似于 D 的中心单 K-代数 B，L 是它的一个极大交换子代数。由于 B 同构于一个矩阵代数 \(\mathbf{M}_{n}(\mathrm{D})\)（VIII, 第 278 页，引理 1），我们有 \([B:K]=n^{2}r^{2}\)，从而 \([L:K]=nr\)（VIII, 第 262 页，命题 3）。

